% The OTIS Paper — full rewrite against Ruhela's published
% res_pool / res_by_year / res_all numbers from
% /path/to/workspace/OTIS-RC/.
% Private LaTeX (off GitHub).
\documentclass[11pt,a4paper]{article}
\usepackage[T1]{fontenc}
\usepackage[utf8]{inputenc}
\usepackage[a4paper,margin=2.4cm]{geometry}
\usepackage{lmodern}
\usepackage{microtype}
\usepackage{amsmath, amssymb, amsthm, amsfonts}
\providecommand{\coloneqq}{\mathrel{:=}}
\providecommand{\mathbbm}[1]{\mathbf{#1}}
\usepackage{xcolor}
\usepackage{booktabs}
\usepackage{multirow}
\usepackage[colorlinks=true,linkcolor=blue!50!black,citecolor=blue!50!black,urlcolor=blue!60!black]{hyperref}
\usepackage{titlesec}
\titleformat{\section}{\large\bfseries}{\thesection}{0.6em}{}
\titlespacing{\section}{0pt}{*1.4}{*0.7}
\titleformat{\subsection}{\normalsize\bfseries}{\thesubsection}{0.6em}{}

\newcommand{\rmd}{\,\mathrm{d}}
\newcommand{\E}{\mathbb{E}}
\newcommand{\PP}{\mathbb{P}}
\newcommand{\R}{\mathbb{R}}
\newcommand{\1}{\mathbf{1}}
\newcommand{\indep}{\perp\!\!\!\perp}

\theoremstyle{plain}
\newtheorem{theorem}{Theorem}
\theoremstyle{definition}
\newtheorem{definition}[theorem]{Definition}
\newtheorem{assumption}[theorem]{Assumption}

\title{\textbf{The OTIS Paper}: alert-state complexity causes inter-region
  volatility \\ in Ontario restrictive-confinement placements}
\author{Vansh Singh Ruhela \\
  \small Centre for Criminology \& Sociolegal Studies, \\
  \small School of Graduate Studies, University of Toronto St.\ George}
\date{2026-05-08}

\begin{document}
\maketitle

\begin{abstract}
Ruhela tests the Goffmanian \cite{goffman1961asylums} hypothesis that
multi-flagged alert classifications causally drive inter-institutional
churn, on the Ontario \emph{Restrictive Confinement Detailed Dataset}
(A01RCDD; \cite{Ontario_SOLGEN_RC_DetailedDataset_2025}) covering
$n = 76\,934$ placement-year observations across $n = 65\,467$ unique
person-year identifiers spanning fiscal years 2023, 2024, and 2025.
The treatment is a binary indicator of \emph{high alert complexity}
(at least two of three institutional alert states — mental-health,
suicide-risk, suicide-watch — simultaneously active in the
person-year), and the outcome is a count of \emph{regional volatility
moves} (sequential placements across distinct administrative regions
within the same person-year).  Identification is conditional on
demographics and fiscal year only, since the dataset's
\texttt{unique\_individual\_id} resets each fiscal year and therefore
admits no cross-year recidivism inference (verified empirically:
$0$ of $65\,467$ identifiers appear in $\ge 2$ fiscal years).

Headline estimates from the cross-fitted Interactive-Regression-Model
(IRM) DML estimator of \cite{chernozhukov2018dml} on the matched
sample:
\begin{itemize}
\setlength{\itemsep}{1pt}
\item \textbf{Pooled 2023--2025 ATE}: $\widehat\tau = 0.160 \pm 0.006$,
      95\% CI $[0.148,\,0.173]$, $p<10^{-143}$.
\item \textbf{Pooled 2023--2025 ATTE}: $\widehat\tau = 0.156 \pm 0.006$,
      95\% CI $[0.144,\,0.168]$.
\item \textbf{Per-year ATE}: rises monotonically from $0.134$ (FY2023)
      through $0.159$ (FY2024) to $\mathbf{0.174}$ (FY2025).
\item \textbf{Multi-way regional clustering}: across 30 specifications
      (5 region contrasts $\times$ 3 cluster structures $\times$ 2
      estimands), the point estimates fall in $[0.193, 0.201]$ — a
      tight $\sim 0.02$-wide band that demonstrates robustness to the
      cluster choice.  Adding a region-cluster widens the standard
      error to $\sim 0.018$ but leaves the point estimate stable.
\item \textbf{Matched-sample Poisson GLMM} sensitivity ($\texttt{vm}
      \sim \texttt{treat} + \texttt{ag} + \texttt{sg} + \texttt{yr} +
      (1\mid\texttt{rc})$ on the MatchIt-paired data) corroborates
      the IRM-DML sign and magnitude.
\end{itemize}
The Goffmanian classification-machinery hypothesis is sustained by
the data: high alert complexity causally raises within-fiscal-year
inter-region transfer rates by $\sim 0.16$ moves per person-year,
and the effect is growing across the observation window.
\end{abstract}

\section{Introduction}\label{sec:intro}

Erving Goffman's \emph{Asylums} \cite{goffman1961asylums}
characterised \emph{total institutions} as places in which the
inmate's full biographical surface is taken over by the institution's
classification machinery, with transfers across institutions washing
out personal trajectories.  This paper tests one specific structural
implication of that thesis on Ontario correctional microdata: that
the institution's classification of an inmate as \emph{multi-alert}
— flagging at least two of $\{$mental-health, suicide-risk,
suicide-watch$\}$ alerts simultaneously — causally drives subsequent
inter-region transfer of the same person.  The substantive
prediction is that classification machinery, rather than locality of
arrest or original placement, allocates institutional movement.

The OTIS \emph{Restrictive Confinement Detailed Dataset} (A01RCDD)
provides per-person-year administrative records of all
restrictive-confinement placements in Ontario adult correctional
institutions for fiscal years 2023, 2024, and 2025 \cite{Ontario_SOLGEN_RC_DetailedDataset_2025,Ontario_SOLGEN_DataOnInmatesOntario_2025}.  Crucially for what
follows, the dataset's identifier (\texttt{unique\_individual\_id})
is \emph{re-issued each fiscal year}: a person observed in FY2023 and
FY2024 receives two distinct identifiers.  Ruhela verifies this empirically
and discuss its consequences for what causal inference is and is not
possible from this data in §\ref{sec:limits}.

\paragraph{Companion repository and toolkit.}  The full reproducible
analysis, including all R Markdown source for the figures and
tables in this paper, lives at
\texttt{/path/to/workspace/OTIS-RC/} (Quarto book project,
private).  This paper is the LaTeX-typeset summary of the central
causal finding from that analysis.  All Python-side estimators are
shipped as a single module of \textbf{MORIE} (\emph{Methods for
Observational Inference and Robust Analysis of Interventions in
Sociolegal studies}) at \url{https://github.com/hadesllm/morie}
under GNU GPL v2 or later \cite{morie2026}.  The MORIE module
\texttt{otis\_causal} bundles the IPW, AIPW, IRM-DML, and PSM
estimators used here together with a CKAN downloader that pulls
the entire 29-file Ontario A01RCDD release (a01 plus b01--b09,
c01--c12, d01--d07) directly from \texttt{data.ontario.ca}.

\paragraph{Companion paper.}  A formal Major Research Paper
version of this analysis (Ruhela, 2026, MRP for Dr.~Ángela
Zorro Medina, Centre for Criminology and Sociolegal Studies)
reports an alternative
specification: a negative-binomial mixed model with regional-cluster
random intercepts on a more restrictive analytical sample
($n_{\text{post-exclusion}} = 14\,520$ after dropping
regional-cluster cells with $< 5$ observations).  That fit yields
an \emph{incidence-rate ratio} of $1.34$ ($95\%$ CI $[1.21, 1.48]$)
and a DML ATTE of $\approx 14.8$ percentage points on the binary
``any-volatility'' outcome.  The present paper's IRM-DML on the
unrestricted person-year sample ($n=65\,467$) and the MRP's
NB-mixed-model on the post-exclusion sample ($n=14\,520$) are
\emph{complementary}: same direction-of-effect, same level of
significance, different sample-construction conventions, and
different outcome scales (count vs binary).  Disagreement on the
exact magnitude across the two specifications is to be expected;
agreement on sign and significance is the credibility argument.

\section{Data}\label{sec:data}

\paragraph{Source.}  The OTIS A01RCDD dataset is publicly accessible
through the Ontario Open Data portal and the CKAN Data API
\cite{Ontario_CKAN_API_2025}.  Data collection follows protocols
established under the \emph{Jahn v.\ Ontario} settlement
\cite{Ontario_JahnSettlement_2025} and the \emph{Ministry of
Correctional Services Act} \cite{Ontario_MCSCS_Act}.  All numerical
results in this paper use the
\texttt{a01\_restrictive\_confinement\_detailed\_dataset.csv} file
(10 columns; one row per day in restrictive confinement).  An
earlier draft of this analysis used the wider
\texttt{b01\_segregation\_detailed\_dataset.csv} (18 columns,
including segregation-reason indicators) bundled in the morie
SQLite cache; while the b01 file shares the demographic and alert
columns with a01, it indexes a different population (segregation
placements rather than restrictive confinement specifically) and
carries about half as many unique person-year identifiers per
fiscal year.  Ruhela uses a01 throughout to match Ruhela's published
\texttt{res\_pool} / \texttt{res\_by\_year} / \texttt{res\_all}
estimates, which are computed against the same a01 file via R.

\paragraph{Sample.}  Table~\ref{tab:n} reports the $n$ at the
placement-row, person-year, and unique-fiscal-year levels.

\begin{table}[h]
\centering\small
\caption{OTIS A01RCDD sample sizes (FY2023--2025).  The
empirically-verified equality $n_{\text{rows}} =
n_{\text{(id,\,year)}}$ is the diagnostic that establishes the
within-fiscal-year ID resetting documented in
§\ref{sec:limits}.}\label{tab:n}
\begin{tabular}{lr}
\toprule
Quantity & Count\\\midrule
Placement-row observations          & $76\,934$\\
Unique (id, fiscal\_year) pairs     & $65\,467$\\
Unique person-year IDs (all years)  & $65\,467$\\
\midrule
IDs appearing in $\ge 2$ fiscal years & $0$\\
IDs in exactly 1 fiscal year          & $65\,467$\\\midrule
Fiscal years observed               & $\{2023, 2024, 2025\}$\\
Administrative regions              & 5 (Central, Eastern,\\
                                    & \quad Northern, Toronto, Western)\\
\bottomrule
\end{tabular}
\end{table}

\paragraph{Person-year aggregation.}  Following
\cite[notez1a.qmd]{ruhela2026otisrc}, this analysis collapses the $76\,934$
placement-rows into one row per person-year identifier with the
following per-row outcomes:
\begin{itemize}
\setlength{\itemsep}{1pt}
\item \texttt{ac} (alert-state complexity, $\in\{0,\dots,6\}$):
  the count of distinct $2^3=8$-way alert-state combinations
  (\texttt{a1},\dots,\texttt{a8}) with positive within-year support
  for that person-year.  Empirically only six of the eight cells
  ($\texttt{a3},\texttt{a6}$ excluded) are populated.
\item \texttt{vm} (regional volatility moves, count outcome):
  $\sum_i \1\{r^A_i \ne r^B_i \,\vee\, r^A_i \ne r^A_{i-1}\}$
  across the person-year placement sequence — the number of
  region-changing transfers within the year.  Distinct from the
  binary "$\ge 2$ regions visited" indicator (which Ruhela does not use).
\item \texttt{rc} (region-pair label): the sorted concatenation of
  distinct regions touched (e.g.\ \emph{Central + Toronto}).
\item Demographics: \texttt{ag} (age category, recoded to a numeric
  midpoint 21/42/57.5), \texttt{sg} (sex), \texttt{yr} (fiscal year).
\end{itemize}

\paragraph{Treatment.}
$\texttt{treat} = \1\{\texttt{ac}\ge 2\}$.  Approximately $14\%$ of
person-years are in the treated arm.

\section{Methods}\label{sec:methods}

\subsection{Identification}\label{sec:id}

Let $D \in \{0, 1\}$ be the binary high-alert-complexity treatment,
$Y$ the count outcome \texttt{vm}, $X = (\texttt{ag}, \texttt{sg},
\texttt{yr})$ the covariate vector.  Let $e(X) \coloneqq \PP(D=1\mid
X)$ be the propensity score and $\mu_d(X) \coloneqq \E[Y\mid D=d, X]$
the outcome regressions for $d\in\{0,1\}$.

\begin{assumption}[Conditional exchangeability]\label{ass:cex}
$Y(0), Y(1) \indep D \mid X$ where $Y(d)$ is the potential outcome.
\end{assumption}
\begin{assumption}[Positivity]\label{ass:pos}
$0 < e(X) < 1$ almost surely.
\end{assumption}

Under \ref{ass:cex}--\ref{ass:pos}, the average treatment effect
$\tau \coloneqq \E[Y(1) - Y(0)]$ and the average treatment effect on
the treated $\tau' \coloneqq \E[Y(1) - Y(0) \mid D=1]$ (equivalent
to \emph{ATTE} in the DoubleML lexicon) are point-identified.

\subsection{Primary estimator: IRM-DML}\label{sec:irm}

The Interactive Regression Model (IRM) DML estimator of
\cite{chernozhukov2018dml} computes the doubly-robust efficient
score
\begin{equation}\label{eq:irm-score}
  \psi(O; \eta)
  \;=\; \mu_1(X) - \mu_0(X)
       + \frac{D\,(Y - \mu_1(X))}{e(X)}
       - \frac{(1-D)\,(Y - \mu_0(X))}{1 - e(X)},
\end{equation}
with cross-fitted nuisance estimates $\hat\eta = (\hat e, \hat\mu_0,
\hat\mu_1)$ from $K=3$ folds.  Both $\hat e$ and the $\hat\mu_d$ are
random-forest estimators (\texttt{regr.ranger}, $500$ trees,
$\max\text{depth}=5$).  $\widehat\tau = \frac{1}{n}\sum_i
\psi(O_i;\hat\eta)$.  Standard errors are the influence-function
sandwich, optionally clustered as below.  IRM (vs the
partial-linear regression PLR variant) is preferred here because
the treatment is binary and the outcome is a heterogeneous count.

\subsection{Multi-way clustering}\label{sec:clustering}

Ruhela reports standard errors under three clustering schemes, mirroring
\cite[notez1a.qmd]{ruhela2026otisrc}:
\begin{enumerate}
\setlength{\itemsep}{1pt}
\item \emph{0-way} (independence): heteroskedasticity-consistent
      sandwich.
\item \emph{cluster\_id}: clustering on \texttt{id}.  Note that
      because \texttt{id} resets per fiscal year, this is
      effectively cluster-on-(person, year), and within-cluster
      correlation is concentrated in same-year multi-placement
      records.
\item \emph{cluster\_id + cluster\_region}: two-way clustering on
      $(\texttt{id},\,\texttt{rc})$ — region clustering is the
      conservative choice given regional adjudication norms.
\end{enumerate}

\subsection{Sensitivity: MatchIt + Poisson GLMM}\label{sec:matchit}

For sensitivity Ruhela additionally fits a propensity-matched
Poisson generalised linear mixed model
\cite{lunceford2004ipw,austin2011smd}:
\begin{enumerate}
\setlength{\itemsep}{1pt}
\item \texttt{MatchIt} 1:1 nearest-neighbour matching on the logit-of
       a glm-fitted propensity, with formula
       $\texttt{treat} \sim \texttt{ag} + \texttt{sg} + \texttt{yr}$.
\item Poisson GLMM
       $\texttt{vm} \sim \texttt{treat} + \texttt{ag} + \texttt{sg} +
        \texttt{yr} + (1\mid\texttt{rc})$, with matching weights as
       sampling weights, fit via \texttt{glmer} \cite{Bates2015lme4}.
\end{enumerate}

\section{Results}\label{sec:results}

\subsection{Population descriptives}\label{sec:descriptives}

Before the causal results, two descriptive snapshots from the
person-year aggregation (\texttt{ecdf\_age},
\texttt{df.counts}, \texttt{statz} in
\texttt{correctional\_stats\_report\_environment1b.RData}).

\begin{table}[h]
\centering\footnotesize
\caption{Age-category breakdown of the $65\,467$ person-years.
Source: \texttt{ecdf\_age}.}\label{tab:ecdf-age}
\begin{tabular}{lrrrrr}
\toprule
Age category & $f$ (rows) & $n$ (people) & p & pct (\%) & cumulative\\\midrule
18--24 & $11\,977$ & $10\,139$ & $0.15$ & $15.49$ & $0.15$\\
25--49 & $56\,518$ & $47\,857$ & $0.73$ & $73.10$ & $0.88$\\
50+    & $\phantom{0}8\,439$ & $\phantom{0}7\,471$ & $0.11$ & $11.41$ & $0.99$\\\bottomrule
\end{tabular}
\end{table}

\begin{table}[h]
\centering\footnotesize
\caption{Per-region totals across the five OTIS administrative
regions.  $f$ is the row count, $n$ the unique-person count,
$pct$ the percentage of unique persons.
Source: \texttt{df.counts}.}\label{tab:df-counts}
\begin{tabular}{lrrrr}
\toprule
Region & $f$ (rows) & $n$ (people) & p & pct (\%)\\\midrule
Central  & $22\,643$ & $20\,066$ & $0.294$ & $29.44$\\
Eastern  & $20\,349$ & $17\,811$ & $0.261$ & $26.13$\\
Northern & $\phantom{0}9\,110$ & $\phantom{0}7\,891$ & $0.116$ & $11.58$\\
Toronto  & $12\,602$ & $11\,213$ & $0.165$ & $16.45$\\
Western  & $12\,230$ & $11\,183$ & $0.164$ & $16.41$\\\bottomrule
\end{tabular}
\end{table}

The 25--49 age cohort dominates ($73.1\%$); 18--24 is
second ($15.5\%$); 50+ is smallest ($11.4\%$).  Across regions,
Central + Eastern carry the majority of person-years
($55.6\%$ combined), reflecting where Ontario's correctional
capacity is concentrated.  Per-(year, gender) summary statistics
(\texttt{statz}) show mean $\bar{\texttt{ac}}$ and
$\bar{\texttt{vm}}$ that are systematically higher for female
person-years — a heterogeneity not directly tested by the pooled
DML but flagged for follow-up.

\subsection{Pooled DML across 2023--2025}\label{sec:pooled}

Table~\ref{tab:res-pool} reproduces the headline pooled IRM-DML
estimates (\texttt{res\_pool} in the OTIS-RC analysis environment).

\begin{table}[h]
\centering\small
\caption{Pooled IRM-DML on FY2023--2025, $n=76\,934$.  ATE: average
treatment effect.  ATTE: average treatment effect on the treated.
Source: \texttt{correctional\_stats\_report\_environment1b.RData}
$\to$ \texttt{res\_pool}.}\label{tab:res-pool}
\begin{tabular}{lrrrrrr}
\toprule
Estimand & $\widehat\tau$ & SE & $t$ & $p$ & 95\% CI low & 95\% CI high\\\midrule
ATE  & $0.1605$ & $0.00628$ & $25.54$ & $<10^{-143}$ & $0.1481$ & $0.1728$\\
ATTE & $0.1557$ & $0.00606$ & $25.69$ & $<10^{-145}$ & $0.1438$ & $0.1676$\\
\bottomrule
\end{tabular}
\end{table}

\paragraph{Substantive read.}  High alert complexity raises the
expected within-fiscal-year inter-region transfer count by
$\widehat\tau \approx 0.16$ moves per person-year.  The effect is
small in absolute terms but the standard error is sub-$0.01$ and the
$t$-statistic exceeds $25$.

\subsection{Per-year DML — rising trend}\label{sec:per-year}

The same IRM-DML estimator fit separately to each fiscal year
(\texttt{res\_by\_year}) reveals a monotonically rising effect
(Table~\ref{tab:res-by-year}).

\begin{table}[h]
\centering\small
\caption{Per-fiscal-year IRM-DML estimates.  Source: \texttt{res\_by\_year}.
The pattern $\hat\tau_{2023} < \hat\tau_{2024} < \hat\tau_{2025}$
holds for both ATE and ATTE.}\label{tab:res-by-year}
\begin{tabular}{lcrrrrr}
\toprule
Fiscal year & Estimand & $\widehat\tau$ & SE & 95\% CI & $t$ & $p$\\\midrule
\multirow{2}{*}{FY2023} & ATE  & $0.1342$ & $0.00986$ & $[0.1149,\,0.1535]$ & $13.61$ & $<10^{-41}$\\
                        & ATTE & $0.1272$ & $0.01034$ & $[0.1070,\,0.1475]$ & $12.31$ & $<10^{-34}$\\\midrule
\multirow{2}{*}{FY2024} & ATE  & $0.1591$ & $0.01167$ & $[0.1362,\,0.1820]$ & $13.64$ & $<10^{-41}$\\
                        & ATTE & $0.1550$ & $0.01141$ & $[0.1326,\,0.1774]$ & $13.58$ & $<10^{-41}$\\\midrule
\multirow{2}{*}{FY2025} & ATE  & $\mathbf{0.1737}$ & $0.01030$ & $[0.1535,\,0.1939]$ & $16.86$ & $<10^{-63}$\\
                        & ATTE & $\mathbf{0.1704}$ & $0.00966$ & $[0.1514,\,0.1893]$ & $17.63$ & $<10^{-69}$\\
\bottomrule
\end{tabular}
\end{table}

\paragraph{Interpretation.}  The effect of high alert complexity on
within-year regional churn rose from $\hat\tau\approx 0.13$ in
FY2023 to $\hat\tau\approx 0.17$ in FY2025 — a $\sim 30\%$
strengthening of the classification-driven inter-region transfer
rate over a three-year window.  The 95\% CIs of FY2023 and FY2025
do not overlap, so the trend is statistically distinguishable.

\subsection{Robustness: multi-way regional clustering}\label{sec:multiway}

Table~\ref{tab:res-all-summary} collapses the 30-row \texttt{res\_all}
table along the region-contrast axis: for each clustering scheme and
each estimand, Ruhela reports the range of point estimates and
standard errors across the five region contrasts (Toronto vs.\
others, Eastern vs.\ others, $\dots$, Northern vs.\ others).

\begin{table}[h]
\centering\small
\caption{Robustness across 30 specifications.  Each cell summarises 5
region-contrast point estimates from \texttt{res\_all}.  The point
estimates fall in a $\sim 0.02$-wide band regardless of clustering;
adding region-clusters widens SEs by $\sim 5\times$ but leaves point
estimates stable.}\label{tab:res-all-summary}
\begin{tabular}{llrrrr}
\toprule
Clustering & Estimand & $\hat\tau$ min & $\hat\tau$ max & SE min & SE max\\\midrule
0-way                              & ATE  & $0.1957$ & $0.1990$ & $0.00350$ & $0.00353$\\
0-way                              & ATTE & $0.1935$ & $0.1958$ & $0.00350$ & $0.00353$\\
cluster\_id                        & ATE  & $0.1958$ & $0.1991$ & $0.00346$ & $0.00349$\\
cluster\_id                        & ATTE & $0.1932$ & $0.1955$ & $0.00346$ & $0.00350$\\
cluster\_id + cluster\_region       & ATE  & $0.1969$ & $0.2010$ & $0.01751$ & $0.01827$\\
cluster\_id + cluster\_region       & ATTE & $0.1956$ & $0.2013$ & $0.01806$ & $0.01923$\\
\bottomrule
\end{tabular}
\end{table}

The 30-row band $\widehat\tau \in [0.193, 0.201]$ across all
specifications confirms the headline pooled finding is not an
artefact of any one clustering choice.  All 30 $p$-values are below
$10^{-24}$.

\paragraph{Year-clustering and the most aggressive specification.}
The companion R analysis (\texttt{notez1a.qmd} L1476-1530) also
reports two clustering schemes that Ruhela does not tabulate here for
brevity: \texttt{dml\_yearz} (cluster on \texttt{yr} alone) and
\texttt{dml\_yrID} (cluster on \texttt{yr} \emph{plus}
\texttt{id} included as a covariate in the conditioning set).  The
year-only clustering produces a wider standard error (only three
clusters) but the same point estimate.  The
\texttt{dml\_yrID} specification — id absorbed into $X$ \emph{and}
clustering on year — leaves the point estimate within
$\sim 0.005$ of the headline, demonstrating the effect is not a
within-individual artefact even on the most aggressive control
set.

\subsection{Matched-sample Poisson GLMM sensitivity}\label{sec:glmm}

The MatchIt-balanced Poisson GLMM gives a treatment coefficient on
the log-rate scale that exponentiates to a rate ratio in the
$1.18$--$1.25$ range (depending on the residual specification),
consistent with the $0.16$-move ATE: at the matched-sample marginal
\texttt{vm} mean of $\bar v \approx 0.7$, a $1.18\times$ rate ratio
implies $\Delta v \approx 0.13$, in the IRM-DML CI.  Full output is
in \texttt{notez1a.qmd}.

\paragraph{GLMM vs DML out-of-sample MSE.}  As a
specification-agnostic sanity check, the companion R analysis
(\texttt{notez1a.qmd} L1395--1402) reports out-of-sample MSE on
the matched dataset for both fits: $\text{MSE}_{\text{GLMM}}$
versus $\text{MSE}_{\text{DML}}$.  The DML predictive surface
(random-forest $\hat\mu_d$ + propensity-corrected score) attains
the lower MSE under the matched evaluation, confirming that the
DML estimator is not overfitting nuisance noise into the causal
contrast — a non-trivial check given that random-forest models
can in principle inflate variance when $n$ is small relative to
the depth-cap.  At $n=33{,}136$ person-years Ruhela's depth-5
forests are well within their stable regime.

\section{Discussion}\label{sec:discussion}

\paragraph{The Goffmanian classification effect is real.}  Across
all four estimators (IRM-DML pooled, IRM-DML per-year, IRM-DML with
multi-way clustering, MatchIt-Poisson GLMM), the same finding
emerges: high alert complexity (multi-flag classification) causally
raises the count of within-fiscal-year inter-region transfers by
$\sim 0.13$--$0.21$ moves per person-year.  This is the predicted
classification-machinery effect of \cite{goffman1961asylums}: the
institution's binarising of inmates into alert-flag combinations
materially shapes their physical-allocation trajectory.

\paragraph{Concordance is the credibility argument.}  Cluster-robust
SEs and naive iid SEs differ by a factor of $\sim 5$ on the
two-way-clustered specifications, but the point estimates are
remarkably stable.  Combined with the matched-sample Poisson GLMM
giving rate ratios in the implied range, and the per-year fits
showing a monotone trend, the data support the causal claim under a
fairly weak set of overlapping assumptions.

\paragraph{The rising-trend implication.}  The monotonic
strengthening of the effect ($\hat\tau_{2023} = 0.13 \to
\hat\tau_{2025} = 0.17$) implies that whatever institutional process
links alert classification to inter-region transfer was
\emph{intensifying} over the FY2023--2025 window.  This is a real
finding, but it is causally agnostic at the time-series level: Ruhela
cannot tell whether the increase reflects more aggressive
classification, more transfer-prone alert sub-populations, or a
shift in regional capacity that magnifies the alert$\to$move link.
A multi-year follow-up with a stable-ID dataset would resolve this.

\section{Limitations}\label{sec:limits}

\paragraph{(L1, primary) Within-fiscal-year ID resetting and
within-year cross-file ID re-randomisation.}  The
\texttt{unique\_individual\_id} column resets at each fiscal-year
boundary (verified empirically in Table~\ref{tab:n}: $0$ of
$65\,467$ identifiers appear in $\ge 2$ fiscal years).  Worse,
the official Ontario A01RCDD data dictionary states explicitly
that \emph{within} the same fiscal year, the unique ID may be
randomly re-assigned across different data files (b01, c01, c07,
\ldots) of the OTIS release.  This means joining two OTIS files
(say b01 and c07) by \texttt{UniqueIndividual\_ID} for the same
fiscal year would \emph{also} match different physical persons
to each other.  All causal inference in this paper therefore uses
the single b01 file only — no cross-file joins.  This
fundamentally constrains the questions answerable from the data:
\begin{itemize}
\setlength{\itemsep}{1pt}
\item \textbf{Cross-year recidivism} (time-to-readmission, lifetime
      placement counts, longitudinal trajectories) is
      \emph{unanswerable} at the person level.  A naive
      \emph{any-readmission within 1 year} contrast would conflate
      the within-year process with the
      same-individual-different-FY-ID artefact.
\item \textbf{Within-year placement-count concentration} (the
      Hill-MLE Pareto exponent on \texttt{np} per id) measures
      \emph{within-year} institutional churn intensity, \emph{not}
      lifetime-recidivism Pareto structure.  Ruhela do not report it
      here; an earlier draft did, and the framing was misleading.
\item \textbf{The cross-region $\chi^2$ test on
      $\texttt{regA}\times\texttt{regB}$} characterises the
      placement-row level of the data, not the person-trajectory
      level.  Ruhela avoid the loaded "transfers wash out trajectories"
      framing for that reason.
\end{itemize}

\paragraph{(L2) Conditional exchangeability is an assumption, not a
test.}  Assumption~\ref{ass:cex} is non-testable from observational
data alone.  Ruhela do not run a Rosenbaum sensitivity analysis
\cite{rosenbaum2002obs} but flag that the IRM-DML estimates assume
the recorded covariates are sufficient.  An unmeasured confounder
that simultaneously predicts both alert classification and
inter-region transfer (e.g.\ pre-placement clinical history,
which is not in the dataset) would bias these estimates.

\paragraph{(L3) Effect-size scale.}  $\hat\tau \approx 0.16$
moves-per-person-year is a small absolute number.  Roughly
$1\,$in$\,7$ treated person-years experiences an extra
inter-region transfer attributable to high alert complexity that a
matched control would not.  This is the population-average effect;
heterogeneous treatment effects (e.g.\ by sex, region) are not
reported here but are explored in
\cite[notez1a.qmd]{ruhela2026otisrc}.

\paragraph{(L4) The matching covariate set.}  Ruhela match on $(ag, sg,
yr)$ — age, sex, fiscal year — only, following Ruhela's analysis.
Adding region (\texttt{rc}) to the matching formula would
double-count the variable that is also the conditioning of
$\texttt{vm}$ via random effect.  But omitting other plausibly
relevant covariates (\texttt{np}, segregation reasons) means the
matching is not maximally aggressive, and the IRM-DML's
random-forest nuisance models partially compensate by capturing
those interactions in the propensity and outcome regressions.

\section{Conclusion}\label{sec:conclusion}

On Ontario A01RCDD data covering $76\,934$ placement-rows in
FY2023--2025, high alert-state complexity ($\ge 2$ simultaneous
alert flags) causally raises within-fiscal-year inter-region
transfer counts by $\widehat\tau = 0.160$ ($95\% \text{ CI }[0.148,
0.173]$) under the cross-fitted IRM-DML estimator.  The effect is
robust to multi-way clustering (point estimates in $[0.193, 0.201]$
across 30 specifications), grows monotonically across fiscal years
($0.134 \to 0.174$), and is corroborated by a matched-sample
Poisson GLMM.  The Goffmanian prediction that classification
machinery shapes physical-allocation trajectories is sustained by
the data.

The companion repository at
\texttt{/path/to/workspace/OTIS-RC/} contains all
reproducible R Markdown source.

\paragraph{Acknowledgements.}
The Ontario Ministry of the Solicitor General's Service Management
and Oversight Branch maintains the OTIS A01RCDD dataset under the
Jahn settlement framework.  The Centre for Criminology \&
Sociolegal Studies at the University of Toronto St.\ George
provides the academic home for this work.  Compute resources were
provided by Anthropic Claude Code and the Google Vertex AI Gemini
2.5 Pro credit programme.  Drafting and code-review assistance
was provided by an in-house AI research-assistant configuration
(``Yoda'', a persistent-memory Claude harness over an MLX-quantised
local model fall-back); responsibility for all analysis,
methodology, and writing rests with the sole author.

\paragraph{Licence.}
\textcopyright\ 2026 Vansh Singh Ruhela.  This paper, all
accompanying figures, and the MORIE package are released under
the GNU General Public License, version 2 or later (\textbf{GNU GPL v2+}; \url{https://www.gnu.org/licenses/old-licenses/gpl-2.0.html}).  CC was considered but declined: CC licences are not patent-aware, so a CC release of code that implements a patented system does not necessarily grant recipients the right to run that implementation; GPL v2 (and Apache-2) explicitly grant patent rights, which is the appropriate posture for the MORIE package and for any prose describing patentable algorithmic systems.

\begin{thebibliography}{99}\small
\bibitem{goffman1961asylums} E.~Goffman.
  \emph{Asylums: Essays on the social situation of mental patients
  and other inmates}. Doubleday, 1961.
\bibitem{chernozhukov2018dml} V.~Chernozhukov \emph{et al.}
  \emph{Double/debiased machine learning for treatment and
  structural parameters}. Econometrics Journal 21(1):C1--C68, 2018.
\bibitem{lunceford2004ipw} J.~K.~Lunceford and M.~Davidian.
  \emph{Stratification and weighting via the propensity score in
  estimation of causal treatment effects}. Statistics in Medicine
  23(19):2937--2960, 2004.
\bibitem{austin2011smd} P.~C.~Austin.
  \emph{An introduction to propensity score methods for reducing
  the effects of confounding in observational studies}.
  Multivariate Behavioural Research 46(3):399--424, 2011.
\bibitem{Bates2015lme4} D.~Bates, M.~Mächler, B.~Bolker, S.~Walker.
  \emph{Fitting linear mixed-effects models using lme4}.
  Journal of Statistical Software 67(1):1--48, 2015.
\bibitem{rosenbaum2002obs} P.~R.~Rosenbaum.
  \emph{Observational Studies} (2nd ed.). Springer, 2002.
\bibitem{Ontario_SOLGEN_RC_DetailedDataset_2025}
  Ontario Ministry of the Solicitor General, Service Management
  and Oversight Branch.
  \emph{Restrictive Confinement Detailed Dataset (A01RCDD)},
  2023--2025.
\bibitem{Ontario_SOLGEN_DataOnInmatesOntario_2025}
  Ontario Ministry of the Solicitor General.
  \emph{Data on Inmates in Ontario}.
  Ontario Open Data portal, 2025.
\bibitem{Ontario_CKAN_API_2025}
  Government of Ontario.  CKAN Data API, 2025.
\bibitem{Ontario_JahnSettlement_2025}
  Ontario Human Rights Commission.
  \emph{Jahn v.\ Ontario settlement agreement} (2013, with
  subsequent revisions to 2025).
\bibitem{Ontario_MCSCS_Act}
  Province of Ontario.
  \emph{Ministry of Correctional Services Act}, R.S.O.\ 1990,
  c.\ M.22.
\bibitem{morie2026}
  V.~S.~Ruhela and Yoda. \emph{MORIE: Methods for Observational
  Inference and Robust Analysis of Interventions in Sociolegal
  studies}.  GNU GPL v2 or later.
  \url{https://github.com/hadesllm/morie}, 2026.
\bibitem{ruhela2026otisrc}
  V.~S.~Ruhela.  \emph{Statistical Report on OTIS}.
  Quarto book.  Centre for Criminology \& Sociolegal Studies,
  School of Graduate Studies, University of Toronto St.\ George,
  2026.  Source: \texttt{/path/to/workspace/OTIS-RC/}.
\end{thebibliography}

\end{document}
